\section{Introduction}
\label{sec:intro}

Valuing an electricity option requires a probability measure under which the
expected discounted payoff is the price. Where a liquid option surface exists,
that measure is recovered from traded instruments: implied volatilities pin down
the diffusion and the forward curve pins down the drift. Neither ingredient is
available in most emerging electricity markets. The Turkish market is a
representative case. Its organised derivative segment, the \emph{Vadeli Elektrik
Piyasası} (VEP) operated by EPİAŞ, lists monthly baseload futures, and
exchange-traded options on power are effectively absent. A practitioner who needs
the value of a call on the hourly spot price at a specific delivery hour
therefore faces two gaps at once: the forward curve is observed only as a coarse
strip of monthly averages, and the risk-neutral distribution around that curve is
not observed at all.

Each gap has been addressed separately in a mature literature. On the
distributional side, regime switching in power prices is well established.
Building on \citet{Hamilton1989}, \citet{Huisman2003} and \citet{Deng2000} argue
that spike behaviour is better represented by a switch between a calm
mean-reverting state and a distinct spike state than by adding jumps to a single
process. \citet{JanczuraWeron2010} compare specifications systematically and find
that independent-spike-regime models with heavy tails dominate simpler
alternatives, consistent with the tail evidence in \citet{Weron2009};
\citet{Deschatre2021} survey the resulting model class. Allowing the switching
probabilities to respond to fundamentals, following \citet{Diebold1994} and
\citet{Filardo1994}, tightens the fit further: \citet{Mount2006} and
\citet{KanamuraOhashi2008} make transition probabilities depend on load and
reserve margin, and \citet{KarakatsaniBunn2008} show that intra-day regime
dynamics carry economic content rather than statistical artefact. Residual
demand, load net of non-dispatchable renewable generation, is the natural
covariate. \citet{Wagner2014} and \citet{Do2021} document its role in price
formation, and the merit-order literature
\citep{Sensfuss2008,Ketterer2014,Clo2015} explains the mechanism: renewable
output displaces the marginal thermal unit and shifts the supply stack. For
Türkiye, \citet{SirinYilmaz2020} quantify the merit-order effect through quantile
regression, while \citet{Ugurlu2018} and \citet{Ertugrul2022} document the
forecasting properties of Turkish spot prices.

On the curve-construction side, turning coarse forward quotes into a granular
curve is equally well studied. \citet{LuciaSchwartz2002} establish the
decomposition of a power price into a deterministic component and a
mean-reverting stochastic residual that underlies most subsequent work;
\citet{FletenLemming2003} pose curve construction as a smoothing problem
constrained to reproduce observed contract prices; \citet{Borak2008} model
forward curve dynamics semi-parametrically; \citet{Kiesel2019} and
\citet{Caldana2017} construct hourly curves for European markets. The common
structure is a smoothness objective subject to arbitrage-free reproduction of the
observed strip.

Pricing on a regime-switching state space is likewise a solved numerical problem.
\citet{BuffingtonElliott2002} and \citet{JobertRogers2006} give the
Markov-modulated framework, \citet{BoyleDraviam2007} and \citet{Huang2011}
develop finite-difference schemes for the coupled systems,
\citet{ChenForsyth2010} apply them to gas storage, and \citet{Zeng2017} compare
finite-difference and Monte Carlo approaches directly. The closest antecedent to
the present study is \citet{Paraschiv2015}, who combine spot and forward
information within a regime-shifting specification. Our departure is to treat the
boundary between the two as an identification question rather than an estimation
one.

That question is what remains open. When a market supplies a forward strip and
nothing else, which parameters of such a model are determined by data, and which
are assumptions the analyst has imported? The distinction is easily obscured. A
model calibrated to reproduce a forward strip to numerical precision looks
calibrated. Where switching parameters and volatilities have been estimated on a
historical sample under the physical measure $\PP$ and are then used for pricing
under $\Q$ without adjustment, the resulting value embeds an untested assumption
that the market charges no premium for regime or volatility risk.
\citet{Janczura2014} constructs such premia explicitly within a regime-switching
framework, \citet{BessembinderLemmon2002} derive equilibrium forward premia from
hedging pressure, and \citet{LongstaffWang2004} and \citet{Algieri2021} measure
them empirically. The premium is a real economic object; it does not vanish
because it was not estimated.

\paragraph{Contribution.}
The separation is made explicit here, and a valuation framework is built around
it. Four contributions follow.

First, a \emph{forward-centred} representation is proposed, in which the spot
price is written as $P_t = F(t) + X_t - \E^{\Q}[X_t]$, with $F$ the
equality-constrained hourly curve and $X$ a two-regime Ornstein--Uhlenbeck
residual with time-varying transition probabilities. Subtracting the residual's
own mean is not a normalisation of convenience: it renders $\E^{\Q}[P_t]=F(t)$ an
algebraic identity holding for every admissible parameter vector, so that the
market-identified first moment is structurally insulated from the unidentified
higher moments. This is established in Proposition~\ref{prop:centering} and
verified numerically across the parameter space explored below, including under
measure changes.

Second, the unidentified inputs are made auditable rather than invisible. Every
parameter carries a provenance label distinguishing calibrated quantities from
inherited historical estimates, derived reconstructions and outright assumptions
(Table~\ref{tab:provenance}), and each is subjected to a controlled sensitivity
experiment in which it alone varies. The risk-neutral drift channel is
implemented as an explicit, switchable perturbation whose default value of zero
reproduces the physical-measure baseline exactly, which makes the zero-premium
assumption a visible choice rather than a silent default.

Third, the framework's principal weakness is shown to lie outside its pricing
layer. Using $5{,}088$ hours of spot prices realised after the valuation date, we
find that the December 2025 forward strip missed the subsequent outcome by a mean
of $-1{,}019$ \TRY, with monthly signed errors from $-24\%$ to $-76\%$, during a
period of exceptionally strong hydrological conditions and record renewable
output. Replaying the curve construction at seven valuation dates from 2022
onward places that outcome in context: it is large but not exceptional, and the
strip's error is negative on average yet unstable in sign. To our knowledge the
predictive performance of the Turkish monthly forward strip has not previously
been measured over a multi-year sample, and the resulting distribution bounds
what forward-anchored valuation can deliver here.

Fourth, a methodological failure and its correction are documented in full. An
initial specification inherited its mean-reversion speed from an autoregressive
fit to the untransformed price level, producing a within-regime half-life of
roughly nineteen years and residual dispersion seven to thirteen times wider than
the realised sample exhibits. Diagnosis required separating within-regime
persistence from the marginal persistence of the switching process, quantities
that differ here by four orders of magnitude. A framework whose sensitivity to a
single inherited coefficient exceeds a factor of four is one whose provenance
discipline is not optional, and the episode is reported on that basis.

Section~\ref{sec:market} describes the market, contract and data;
Section~\ref{sec:curve} develops the hourly curve; Section~\ref{sec:residual}
specifies the residual dynamics and the centring identity;
Section~\ref{sec:valuation} derives the moment system and pricing equation;
Section~\ref{sec:provenance} sets out parameter provenance; and
Section~\ref{sec:results} reports calibration quality, the option surface, the
sensitivity battery and the realised evaluation.
